\documentclass{article}

\usepackage{amsmath,amssymb}
\usepackage{xurl}
\usepackage[hidelinks]{hyperref}
\setlength{\emergencystretch}{2em}

% Shared commands for notes in docs/paper-gaps/.

\DeclareUnicodeCharacter{2080}{\ifmmode{}_{0}\else\textsubscript{0}\fi}
\DeclareUnicodeCharacter{2081}{\ifmmode{}_{1}\else\textsubscript{1}\fi}
\DeclareUnicodeCharacter{2082}{\ifmmode{}_{2}\else\textsubscript{2}\fi}
\DeclareUnicodeCharacter{2099}{\ifmmode{}_{n}\else\textsubscript{n}\fi}

\newcommand{\Lean}{\textsc{Lean}}
\ExplSyntaxOn
\NewDocumentCommand{\leanid}{m}{
  \ifmmode
    \textsf{\detokenize{#1}}
  \else
    \group_begin:
    \urlstyle{sf}
    \tl_set:Nn \l_tmpa_tl {#1}
    \tl_replace_all:Nnn \l_tmpa_tl {\_} {_}
    \exp_args:NV \path \l_tmpa_tl
    \group_end:
  \fi
}
\ExplSyntaxOff
\providecommand{\tr}{\operatorname{tr}}
\newcommand{\tnleanIssueBase}{https://github.com/LionSR/TNLean/issues/}
\newcommand{\tnleanPrBase}{https://github.com/LionSR/TNLean/pull/}
\newcommand{\tnleanDocBase}{https://sirui-lu.com/TNLean/docs/}
\newcommand{\ghissue}[1]{\href{\tnleanIssueBase#1}{GitHub issue \##1}}
\newcommand{\ghpr}[1]{\href{\tnleanPrBase#1}{PR \##1}}
\newcommand{\doclink}[1]{\href{\tnleanDocBase#1.html}{\path{#1}}}

% Dirac notation and the identity operator, as in blueprint/src/macros/common.tex.
% \providecommand keeps note-local definitions authoritative.
\usepackage{dsfont}
\providecommand{\ket}[1]{|#1\rangle}
\providecommand{\bra}[1]{\langle #1|}
\providecommand{\braket}[2]{\langle #1 | #2 \rangle}
\providecommand{\ketbra}[2]{|#1\rangle\!\langle #2|}
\providecommand{\Id}{\mathds{1}}
\providecommand{\one}{\mathds{1}}
\providecommand{\C}{\mathbb{C}}
\providecommand{\MN}[1]{M_{#1}(\C)}
\providecommand{\MD}{M_D(\C)}

% Verdict marker: \gapnote{<kind>}{<status>}. Kind is the policy
% classification (clarification, local-correction, scope-restriction,
% unfaithful, false-source, open-gap); status is open, wip (elimination
% underway), resolved, or historical. Machine-read by the site index;
% prints a verdict line.
\newcommand{\gapnote}[2]{\par\noindent\textbf{Verdict:} #1 (#2).\par}


% Fallbacks keep this note compilable when a different command.tex is found first.
\providecommand{\leanid}[1]{\textsf{\detokenize{#1}}}
\providecommand{\gapnote}[2]{\par\noindent\textbf{Verdict:} #1 (#2).\par}

\title{The Uniform Repair Induction from the Initial Isolated Stars}
\author{}
\date{2026-10-15}

\begin{document}
\maketitle
\gapnote{scope-restriction}{open}

\section{The source lemma and the formal restriction}

Lemma~11.4 of the September~24, 2026 two-dimensional area-law
manuscript~\cite{OpenAI2026AreaLaw} (\path{10-geometry.tex}, lines~471--485)
asserts, for a sufficiently small fixed isolation constant $a_0$ and the
corresponding stopping scale $S_*$, that every active mark $(v,S)$ of every
generation of the repair has a terminal-square-free isolated star
$v+[-a_0S,a_0S]^2$, that distinct marks satisfy the separation inequality,
that the simultaneous removal squares are disjoint, that same-label contacts
of replacements occur only at marks of the next generation, that descendants
of $(v,S)$ stay within distance $2d(S)$ of $v$, and that every branch
terminates. Its proof (lines~487--543) is an induction whose base case is
Lemma~11.3 (lines~332--349), the isolated stars of the initial mesh.

The formal theorem proves the full conclusion of Lemma~11.4 with the
conclusion of Lemma~11.3 for the initial repair stage as an explicit
hypothesis. No other hypothesis is added: the isolation constant is the
smaller of the initial constant and $1/4$ at every generation, and the
stopping scale and the lower layer index are taken large, as in the source.

\section{Elimination}

Lemma~11.3 is proved separately for the jointly chosen initial stage.
Applying the conditional theorem to that proof gives Lemma~11.4 with the
source's hypotheses. This note is then obsolete and is deleted together with
the scope marker.

\section{Lean declarations}

\begin{itemize}
    \item \leanid{TNLean.PEPS.AreaLaw.Geometry.InitialIsolatedStarsStatement}
    \item \leanid{TNLean.PEPS.AreaLaw.Geometry.UniformRepairInductionStatement}
    \item \leanid{TNLean.PEPS.AreaLaw.Geometry.uniformRepairInductionStatement_of_initialIsolatedStars}
\end{itemize}

\bibliographystyle{alpha}
\bibliography{references}

\end{document}
